% from ega4-10-fr.tex
\oldpage[IV-3]{96}
\begin{proposition}[10.1.2]
\label{IV.10.1.2-fr}
Soient $X$ un espace topologique, $X_0$ une partie de $X$. Les conditions
suivantes sont équivalentes~:
\begin{enumerate}[label=\emph{\alph*)}]
  \item Pour toute partie localement fermée $Z\neq\varnothing$ de $X$, on a
  $Z\cap X_0\neq\varnothing$.
  \item[\emph{a$'$)}] Pour toute partie fermée $Z$ de $X$, on a
  $Z=\overline{Z\cap X_0}$.
  \item Pour toute partie $Z\neq\varnothing$ de $X$ localement
  quasi-constructible, on a $Z\cap X_0\neq\varnothing$.
  \item[\emph{b$'$)}] Pour toute partie localement quasi-constructible $Z$ de
  $X$, on a $Z\subset\overline{Z\cap X_0}$ (autrement dit $Z\cap X_0$ est
  \emph{dense} dans $Z$).
  \item L'application $U\mapsto U\cap X_0$ de $\mathfrak O(X)$ dans
  $\mathfrak O(X_0)$ est injective (donc bijective).
  \item[\emph{c$'$)}] L'application $F\mapsto F\cap X_0$ de $\mathfrak F(X)$
  dans $\mathfrak F(X_0)$ est injective (donc bijective).
  \item[\emph{c$''$)}] L'application $Z\mapsto Z\cap X_0$ de
  $\mathfrak{Qc}(X)$ dans $\mathfrak{Qc}(X_0)$ est injective (donc bijective).
  \item[\emph{c$'''$)}] L'application $Z\mapsto Z\cap X_0$ de
  $\mathfrak{Lqc}(X)$ dans $\mathfrak{Lqc}(X_0)$ est injective.
\end{enumerate}
En outre, lorsque ces conditions sont vérifiées, l'application
$Z\mapsto Z\cap X_0$ de $\mathfrak{Lqc}(X)$ dans $\mathfrak{Lqc}(X_0)$ est
bijective.
\end{proposition}

Notons que les assertions de surjectivité dans \emph{c)} et \emph{c$'$)} sont
triviales~: elles entraînent que toute partie localement fermée dans $X_0$ est
trace sur $X_0$ d'une partie localement fermée dans $X$, donc l'application
$\mathfrak{Qc}(X)\to\mathfrak{Qc}(X_0)$ définie dans \emph{c$''$)} est aussi
surjective.

Nous prouverons les implications
\[
c''')\Rightarrow c'')\Rightarrow c')\Rightarrow c)\Rightarrow b)
\Rightarrow b')\Rightarrow a')\Rightarrow a)\Rightarrow c''').
\]
Les trois premières sont triviales. Pour voir que \emph{c)} implique
\emph{b)}, notons d'abord que \emph{c)} entraîne que $X_0$ est \emph{dense}
dans $X$. Remplaçant $X$ et $X_0$ par un ouvert convenable $U$ de $X$ et par
$U\cap X_0$ respectivement, on peut donc se borner au cas où $Z$ est
localement fermée dans $X$, donc $Z=V\cap\complement W$, où $V$ et $W$ sont
ouverts dans $X$~; l'hypothèse $Z\neq\varnothing$ signifie que
$V\not\subset W$, ou encore $V\cup W\neq W$. En vertu de \emph{c)}, on a donc
$(V\cup W)\cap X_0\neq W\cap X_0$, donc $V\cap X_0\not\subset W$, et par
suite $(V\cap\complement W)\cap X_0\neq\varnothing$.

Pour voir que \emph{b)} entraîne \emph{b$'$)}, il suffit d'appliquer
\emph{b)} à $Z\cap U$, où $U$ est un voisinage ouvert arbitraire d'un point de
$Z$. Comme $\overline{Z\cap X_0}\subset\overline{Z}$, il est trivial
que \emph{b$'$)} implique \emph{a$'$)}. Pour montrer que \emph{a$'$)} entraîne
\emph{a)}, remarquons que si $Z$ est localement fermée dans $X$, on peut
écrire $Z=F-F'$ où $F$ et $F'$ sont fermés dans $X$ et $F'\subset F$~; d'où
$Z\cap X_0=(F\cap X_0)-(F'\cap X_0)$. Si l'on avait
$Z\cap X_0=\varnothing$, on en déduirait $F\cap X_0=F'\cap X_0$, donc
$F=F'$ en vertu de \emph{a$'$)}, c'est-à-dire $Z=\varnothing$.

Pour voir que \emph{a)} entraîne \emph{c$'''$)}, il suffit de montrer que si
$Z\neq\varnothing$ est localement quasi-constructible, on a
$Z\cap X_0\neq\varnothing$~: en effet, la relation
$Z'\cap X_0=Z''\cap X_0$ est équivalente à
$((Z'\cup Z'')-(Z'\cap Z''))\cap X_0=\varnothing$. Autrement dit, il suffit de
prouver que \emph{a)} entraîne \emph{b)}~; en outre, remplaçant $X$ et $X_0$
par un ouvert $U$ de $X$ et par $U\cap X_0$ respectivement, on est bien
ramené au cas où $Z$ est localement fermée dans $X$, d'où la conclusion.

Reste à montrer que l'application
$\mathfrak{Lqc}(X)\to\mathfrak{Lqc}(X_0)$ est surjective. Soit donc $Z_0$ une
partie localement quasi-constructible de $X_0$~: il y a un recouvrement
$(V_\alpha)$ de $X_0$ par des ouverts de $X_0$ tel que $Z_0\cap V_\alpha$ soit
quasi-constructible dans $V_\alpha$ (et
